% One-time imported draft. Editable native source; no Markdown regeneration.
\section{Trace an oscillator without mistaking repetition for
inactivity}\label{008-life-generations-p09-trace-an-oscillator-without-mistaking-repetition-for-inactivity}

The default live set starts at (2,3), (3,3), and (4,3). The middle cell
has two neighbors and survives. The two ends have one each and die. The
positions (3,2) and (3,4) each have three neighbors and are born, giving
a vertical line in \texttt{board\_b}.

Applying the same synchronous rule to that vertical state returns a
horizontal line in \texttt{board\_a}. The visible generation-two image
matches generation zero, but \texttt{board\_b} still retains generation
one. The ready checkpoints and buffer contents distinguish real
evolution from an implementation that simply renders the seed twice.

A 2 by 2 live block is a different example. Each live cell has three
neighbors and survives, while the surrounding dead cells lack the
three-neighbor configuration needed for birth. It is stable under this
finite model when placed away from an edge. Repeated images therefore
need interpretation: some repeat every step, others alternate.
